\documentclass[letterpaper,fleqn]{article}
\usepackage[spanish,es-noshorthands]{babel}
\usepackage[utf8]{inputenc} 
\usepackage[height=9.5in,width=7in]{geometry}
\usepackage{mathexam}
\usepackage{amsmath}
\usepackage{graphicx}

\ExamClass{\includegraphics[height=16pt]{Images/logo-sed.png} Álgebra $8^{\circ}$}
\ExamName{Prueba Bimestral III}
\ExamHead{\includegraphics[height=16pt]{Images/logo-colegio.png} IEDAB}
\newcommand{\LineaNombre}{%
\par
\vspace{\baselineskip}
Nombre:\hrulefill \; Curso: \underline{\hspace*{48pt}} \; Form: \underline{A} \; Fecha: \underline{\hspace*{2.5cm}} \relax
\par}
\let\ds\displaystyle

\begin{document}
\ExamInstrBox{
Respuesta sin justificar mediante procedimiento no será tenida en cuenta en la calificación. Escriba sus respuestas en el espacio indicado. Tiene 70 minutos para contestar esta prueba. Cada punto de esta prueba vale 0.8 puntos para un total de 4.0 puntos y 1 punto por presentar la prueba.}
\LineaNombre
\begin{enumerate}
 \item Exprese usando lenguaje algebraico las siguientes expresiones:
 \begin{enumerate}
  \item El cuadrado de la diferencia de dos números. \hspace{1cm} Resp: \hrulefill 
  \item La suma del triple de un número y el cuádruple de otro \hspace{1cm} Resp: \hrulefill
  \item La raíz cuadrada del producto de dos números. \hspace{1cm} Resp: \hrulefill
  \item La mitad de la suma de dos números. \hspace{1cm} Resp: \hrulefill
  \item El cubo de la suma del doble de un número y el triple de otro \hspace{1cm} Resp: \hrulefill
 \end{enumerate}
\item Exprese en lenguaje verbal las siguientes expresiones algebraicas
\begin{enumerate}
 \item $(a-b)^{2}$:\hspace{1cm} Resp: \hrulefill
 \item $3(x+y)$: \hspace{1cm} Resp: \hrulefill
 \item $\dfrac{x-y}{3}$:\hspace{1cm} Resp: \hrulefill
 \item $\sqrt{2x+5y}$:\hspace{1cm} Resp: \hrulefill
\end{enumerate}
\item Resuelva los siguientes problemas, planteando la ecuación correspondiente y solucionándola:
\begin{enumerate}
 \item  Entre Luis y Antonio reúnen 840 euros. Sabiendo que Antonio tiene 125 euros más que Luis, calcular los euros que tiene cada uno. \noanswer
  \item Repartir 300 euros entre tres personas de modo que la segunda reciba 16 euros más que la primera y la tercera 28 euros más que la segunda. \noanswer
\end{enumerate}
\newpage
\item Dados los polinomios $P=3x^{2}-4x+5$ \hspace{24pt} y \hspace{24pt} $Q=2x^{2}+6x-8$ halle:
\begin{enumerate}
 \item $P+Q=$\noanswer
 \item $P-Q=$\noanswer
 \item $4P=$\noanswer
 \item $P\cdot Q=$\noanswer
\end{enumerate}
\item Desarrolle usando la propiedad distributiva de la multiplicación respecto a la suma o cuando pueda use los productos notables
\begin{enumerate}
 \item $8a^{3}b^{4}(3ab-2ab^{2}+4a^{2}b^{2})=$\noanswer
 \item $(6x+7)(3x-10)=$\noanswer
 \item $(3t+7)^{2}=$\noanswer
 \item $(4x-2)^{3}=$\noanswer
 \item $(5x-2a)(5x+2a)=$\noanswer
\end{enumerate}
 \end{enumerate}

\end{document}
